% Source: https://artofproblemsolving.com/community/c3h1666962_concept_regarding_2018_state_target__6
% Stable identity: aops-topic-1666962-q01
\textbf{(a)} The line $y=mx$ with $m>0$ contains the ray from $O$ into the first quadrant making the angle $\alpha$ with the positive $x$-axis, where $\tan\alpha=m$ and $0^\circ<\alpha<90^\circ$. Similarly $y=nx$ contains the ray at angle $\beta$ with $\tan\beta=n$, and $\beta<\alpha$ because $n<m$. The angle between these two rays is $\alpha-\beta<90^\circ$; the other angle between the lines (between the first ray and the opposite of the second) is its supplement.

The line $y=x$ makes the angle $45^\circ$ with the positive $x$-axis. It can bisect only the angle between the two first-quadrant rays: the bisector of the other angle points in the direction $\dfrac{\alpha+\beta+180^\circ}{2}$, which lies strictly between $90^\circ$ and $180^\circ$, so it is not the line $y=x$ (whose directions are $45^\circ$ and $225^\circ$). Hence $y=x$ bisects the angle from $\beta$ to $\alpha$, which means $\dfrac{\alpha+\beta}{2}=45^\circ$. Writing $\theta$ for each half,
\[
\alpha=45^\circ+\theta,\qquad \beta=45^\circ-\theta,\qquad \alpha+\beta=90^\circ .
\]

\begin{asy}
usepackage("amsmath");
size(5cm);
defaultpen(fontsize(10pt));
pair O=(0,0);
draw((-0.3,0)--(4.4,0),gray,Arrow(TeXHead));draw((0,-0.3)--(0,4.4),gray,Arrow(TeXHead));
label("$x$",(4.4,0),E,gray);label("$y$",(0,4.4),N,gray);
draw(O--(4/3,4),linewidth(0.9));draw(O--(4,4/3),linewidth(0.9));
draw(O--(3.7,3.7),blue+dashed);
real al=aTan(3),be=aTan(1/3);
draw(arc(O,1.4,45,al),red);draw(arc(O,1.4,be,45),red);
label("$\theta$",1.75*dir((45+al)/2),red);label("$\theta$",1.75*dir((45+be)/2),red);
label("$y=mx$",(4/3,4),N);label("$y=nx$",(4,4/3),E);label("$y=x$",(3.7,3.7),NE,blue);
label("$O$",O,SW);
label("\mbox{}",(min(currentpicture,user=true)+max(currentpicture,user=true))/2,black); // keeps the text after the figure black when LaTeX typesets the labels inline
\end{asy}

Therefore
\[
n=\tan\beta=\tan(90^\circ-\alpha)=\dfrac{1}{\tan\alpha}=\dfrac{1}{m},\qquad\text{so } mn=1 .
\]
(Equivalently: reflection in $y=x$ swaps the two lines, and it sends the line $y=mx$ to $x=my$, the line of slope $\dfrac{1}{m}$. The figure uses $m=3$, $n=\dfrac{1}{3}$.) $\blacksquare$

\textbf{(b)} By (a), $mn=1$, so
\[
(m-n)^2=(m+n)^2-4mn=4\cdot65-4=256 .
\]
Since $m>n$, $m-n=\boxed{16}$. Indeed $m=\sqrt{65}+8$ and $n=\sqrt{65}-8$, with product $65-64=1$.
